---
title: "Fourier Analysis on Finite Groups"
author: "Eugen Nesbakken"
date: "2026-02-23"
description: "An introduction to Fourier analysis on finite abelian groups and its connections to representation theory."
category: mathematics
tags: [algebra, fourier-analysis, representation-theory]
---

\section{Introduction}

Fourier analysis on finite groups generalises the classical discrete Fourier transform to an arbitrary finite abelian group $G$.
The central object is the \emph{Pontryagin dual} $\hat{G}$, the group of characters of $G$.

\section{Characters and the Dual Group}

\begin{definition}[Character]
A \textbf{character} of a finite abelian group $G$ is a group homomorphism
\[
  \chi : G \to \mathbb{C}^{\times}.
\]
The set of all characters forms the \textbf{Pontryagin dual} $\hat{G}$, which is itself a finite abelian group under pointwise multiplication.
\end{definition}

For $G = \mathbb{Z}/n\mathbb{Z}$, the characters are $\chi_k(m) = e^{2\pi i km/n}$ for $k = 0, 1, \ldots, n-1$, and $\hat{G} \cong G$.

\section{The Fourier Transform}

\begin{definition}[Fourier Transform]
Let $f : G \to \mathbb{C}$.
The \textbf{Fourier transform} of $f$ is the function $\hat{f} : \hat{G} \to \mathbb{C}$ defined by
\[
  \hat{f}(\chi) = \sum_{g \in G} f(g)\,\overline{\chi(g)}.
\]
\end{definition}

The Fourier transform is inverted by the \textbf{Fourier inversion formula}:
\[
  f(g) = \frac{1}{|G|} \sum_{\chi \in \hat{G}} \hat{f}(\chi)\,\chi(g).
\]

\subsection{Orthogonality of Characters}

The characters satisfy the orthogonality relations
\[
  \sum_{g \in G} \chi(g)\,\overline{\psi(g)} = |G|\,\delta_{\chi,\psi},
\]
which make $\{|G|^{-1/2}\chi\}_{\chi \in \hat{G}}$ an orthonormal basis for $L^2(G) = \mathbb{C}^G$.

\section{Convolution and Plancherel}

The \textbf{convolution} of $f, h : G \to \mathbb{C}$ is
\[
  (f * h)(g) = \sum_{x \in G} f(x)\,h(g - x).
\]
Under the Fourier transform, convolution becomes pointwise multiplication: $\widehat{f * h}(\chi) = \hat{f}(\chi)\,\hat{h}(\chi)$.

\begin{theorem}[Plancherel]
There is an isometry $L^2(G) \cong L^2(\hat{G})$ given by
\[
  \sum_{g \in G} |f(g)|^2 = \frac{1}{|G|} \sum_{\chi \in \hat{G}} |\hat{f}(\chi)|^2.
\]
\end{theorem}

\begin{proof}
Apply the Fourier inversion formula and the orthogonality of characters.
\end{proof}

\section{Connection to Representation Theory}

For a non-abelian finite group $G$, characters of one-dimensional representations are replaced by \emph{matrix coefficients} of irreducible representations.
The Peter--Weyl theorem (in its finite-group form) states that the functions
\[
  g \mapsto \sqrt{\dim V}\,\rho_{ij}(g), \qquad \rho : G \to \GL(V) \text{ irreducible},
\]
form an orthonormal basis for $L^2(G)$, recovering and extending the abelian theory.
